% =====================================================================
% DRAFT: METHODS FORMALIZATION
% Target: AAAI_27_Rodent/AnonymousSubmission2027.tex
% Replaces: "Feature Engineering Framework" (lines 173--191) and adds
%           a feature table and an algorithm box.
% Revision-plan items: A6.3 (Welch), A6.4 (MMD), A6.5 (Hjorth),
%           A6.6 (feature table), A6.7 (algorithm box), plus the CFC
%           slot demanded by section 6.1 / prerequisite P1.
%
% TRACK: this file is TRACK A (writing only, zero new experiments) with
% the single exception noted in the CFC subsection, which is Track B if
% and only if no coupling quantity is currently computed.
%
% ---------------------------------------------------------------------
% !! TWO PREAMBLE LINES ARE REQUIRED AND ARE NOT PRESENT YET !!
%
% The algorithm box below needs these two lines added to the preamble of
% AnonymousSubmission2027.tex, immediately after \frenchspacing:
%
%     \usepackage{algorithm}
%     \usepackage{algorithmic}
%
% They are NOT on the AAAI-27 forbidden list, they are the packages the
% AAAI author kit itself recommends for algorithms, and the companion
% paper at AAAI_27/AnonymousSubmission2027.tex already loads exactly
% these two (its preamble comment reads "These are recommended to typeset
% algorithms but not required"). Both .sty files are present in the local
% TeX Live installation. Nothing else new is required: everything below
% uses amsmath, amssymb, booktabs and natbib, all already loaded.
%
% If the authors decide not to add those two packages, delete the
% algorithm float and nothing else in this file breaks.
% ---------------------------------------------------------------------
%
% NEW BIB KEYS USED HERE (all in new_refs.bib, all verified):
%   hjorth1970, welch1967, tort2010, canolty2010
% EXISTING BIB KEYS USED HERE:
%   e18090272 (Aboalayon et al.), Chen_2016
%
% CONVENTION USED THROUGHOUT THIS FILE:
%   Any quantity that is not already stated in the manuscript is left as
%   a symbol and flagged with a
%       % REQUIRES_EXPERIMENT: <what to run / read from the code>
%   comment on the preceding line. There are no placeholder numbers.
%   Most of these are not experiments in the usual sense -- they are
%   values that must be transcribed from the feature-extraction code.
%   They are flagged the same way because writing a guess would be
%   fabrication either way.
% =====================================================================


\subsection{Notation}
\label{subsec:notation}

% NOTE TO THE AUTHORS: if the separate "Problem setting" paragraph from
% item A6.1 of the revision plan is also adopted, merge it with this
% paragraph rather than keeping both -- they define the same symbols.
Let $r$ denote a continuous single-channel EEG recording sampled at
$f_s = 500$\,Hz. We segment $r$ into non-overlapping epochs of duration
$T = 10$\,s, so that each epoch is a vector
$x = (x_1,\dots,x_N) \in \mathbb{R}^{N}$ with $N = T f_s = 5000$ samples.
A feature map $\phi : \mathbb{R}^{N} \to \mathbb{R}^{d}$ produces the
descriptor vector defined in this section, and a classifier
$f : \mathbb{R}^{d} \to \Delta^{2}$ returns a distribution over the three
vigilance states $\mathcal{Y} = \{\textsc{wake}, \textsc{sws},
\textsc{rem}\}$. Table~\ref{tab:features} enumerates the components of
$\phi$ and Algorithm~\ref{alg:pipeline} gives the end-to-end procedure.


\subsection{Spectral Features}
\label{subsec:spectral}

Each epoch is converted to a one-sided power spectral density (PSD)
estimate by Welch's method \citep{welch1967}. The epoch is divided into
$K$ segments of $L$ samples overlapping by $D$ samples, written
$x^{(k)}_n = x_{\,n + (k-1)(L-D)}$ for $n = 0, \dots, L-1$. Each segment is
multiplied by a window $w \in \mathbb{R}^{L}$ with energy
$U = \sum_{n} w_n^{2}$, and the resulting modified periodograms are
averaged:
%
\begin{equation}
\label{eq:welch}
\hat{S}(f) \;=\;
\frac{1}{K f_s U}
\sum_{k=1}^{K}
\left|
\sum_{n=0}^{L-1} w_n\, x^{(k)}_n\,
e^{-\mathrm{j} 2\pi f n / f_s}
\right|^{2}.
\end{equation}
%
The estimator is fully determined by the window function $w$, the segment
length $L$, the overlap $D$, the transform length $N_{\mathrm{FFT}}$, the
per-segment detrending rule, and the output scaling; the achievable
frequency resolution is $\Delta f = f_s / L$.

% REQUIRES_EXPERIMENT: transcribe from the feature-extraction code the
% exact scipy.signal.welch arguments -- window, nperseg (L), noverlap (D),
% nfft, detrend, scaling ('density' vs 'spectrum') -- and report the
% resulting Delta f. Do not guess: with f_s = 500 Hz, any segment shorter
% than 2 s cannot resolve the lower edge of the delta band at all, so this
% parameter materially determines whether the headline delta feature is
% well defined.
\noindent
% Fill the four bracketed slots below from the code, then delete this note.
We use a [window] window with $L = [\;]$ samples and $D = [\;]$ samples of
overlap, giving $\Delta f = [\;]$\,Hz.

Five canonical bands are used:
$\delta = [0.5, 4)$\,Hz,
$\theta = [4, 8)$\,Hz,
$\alpha = [8, 12)$\,Hz,
$\beta = [12, 30)$\,Hz and
$\gamma = [30, 100)$\,Hz.
% REQUIRES_EXPERIMENT: confirm the half-open band-edge convention above
% against the code. The manuscript writes the bands as closed intervals
% (0.5--4, 4--8, ...), which leaves the assignment of a bin falling
% exactly on an edge undefined. State whichever convention the code uses;
% the interval notation above is a proposal, not a measurement.
Writing $B_b \subset [0, f_s/2]$ for the frequency support of band $b$,
we compute for each band its absolute power, its relative power, and its
peak frequency:
%
\begin{align}
P_b       &= \Delta f \!\!\sum_{f \in B_b}\!\! \hat{S}(f),
\label{eq:bandpower}\\
\tilde{P}_b &= P_b \Big/ \sum_{b' \in \mathcal{B}} P_{b'},
\label{eq:relpower}\\
f_b^{\star} &= \operatorname*{arg\,max}_{f \in B_b} \hat{S}(f).
\label{eq:peakfreq}
\end{align}
%
% REQUIRES_EXPERIMENT: three conventions to transcribe here, each of which
% changes the numbers: (i) whether band power is the integral of
% Eq. (2) or the mean PSD over the band; (ii) whether the denominator of
% Eq. (3) is the sum over the five bands, as written, or total power over
% the full 0.5--100 Hz support; (iii) whether a log transform was applied
% to P_b before it entered the design matrix.
Finally, normalizing the PSD to a distribution
$p(f) = \hat{S}(f) / \sum_{f' \in F} \hat{S}(f')$ over the analysis support
$F$, we take the normalized spectral entropy
%
\begin{equation}
\label{eq:spectralentropy}
H(x) \;=\; -\frac{1}{\log |F|} \sum_{f \in F} p(f) \log p(f)
\;\in\; [0,1],
\end{equation}
%
which is $0$ for a pure tone and $1$ for a flat spectrum and therefore
indexes how broadband an epoch is.
% REQUIRES_EXPERIMENT: state the support F (full spectrum, or per band)
% and the logarithm base. The design matrix underlying Figure 3 contains a
% single column named "entropy", whereas the current Methods prose says
% spectral entropy was computed "for each frequency band". Those two
% statements cannot both be right; resolve against the code before
% filling the dimensionality of this row in Table 1.


\subsection{Hjorth Descriptors}
\label{subsec:hjorth}

% NOTE: this subsection is new. Hjorth activity, mobility and complexity
% are the two highest-ranked features in the paper's own Figure 3 and are
% currently introduced for the first time in the Results section, with no
% definition and no citation. "activity" is never mentioned in the text at
% all. Item D13 of the revision plan.
Alongside the frequency-domain descriptors we compute the three time-domain
descriptors introduced by \citet{hjorth1970}. Let $x'$ and $x''$ denote the
first and second derivatives of the epoch. \emph{Activity} is the signal
variance, \emph{mobility} is the ratio of the standard deviation of the
derivative to that of the signal, and \emph{complexity} is the mobility of
the derivative relative to the mobility of the signal:
%
\begin{align}
A(x) &= \operatorname{Var}(x), \label{eq:activity}\\
M(x) &= \sqrt{\operatorname{Var}(x') \big/ \operatorname{Var}(x)},
\label{eq:mobility}\\
C(x) &= M(x') \big/ M(x). \label{eq:complexity}
\end{align}
%
These are not independent of the spectral features defined above.
% NOTE: do not replace this with \ref to a subsection label. The template
% sets secnumdepth to 0, so subsections carry no number and \ref to a
% subsection label prints a stale counter value, not a section number.
Writing $m_k = \int \omega^{k} S(\omega)\,d\omega$
for the $k$-th moment of the power spectrum, one has $A = m_0$,
$M = \sqrt{m_2/m_0}$ and $C = \sqrt{m_4 m_0 / m_2^{2}}$: mobility is a
normalized mean frequency and complexity is a normalized spectral
bandwidth. Both are therefore functions of the same spectrum that produces
$P_b$ and $\tilde{P}_b$, and are strongly collinear with the band-power
features by construction. We note this explicitly because gain-based
importance attribution distributes credit arbitrarily among collinear
predictors, which bears directly on how the rankings reported later should
be read.

% REQUIRES_EXPERIMENT: state (i) whether x' and x'' are computed as
% first differences (x'_n = x_n - x_{n-1}) or as scaled derivatives
% (divided by 1/f_s), and (ii) whether they are computed on the raw or on
% a filtered signal. Point (i) is not cosmetic: at f_s = 500 Hz the two
% conventions differ by a factor of 500 in mobility. Complexity, being a
% ratio of two mobilities, is invariant to the choice; mobility is not.
% Since mobility is the second-ranked feature in Figure 3, the convention
% must be stated for the value to be reproducible.


\subsection{Maximum--Minimum Distance}
\label{subsec:mmd}

% ATTRIBUTION FIX (item A6.4 / D12): the words "custom" (line 189) and
% "novel" (line 201) must be deleted from the manuscript. The feature is
% attributed to the reference cited in the very same sentence. The
% abstract only lists MMD as a feature used, never as a new one, so this
% correction is abstract-safe.
The Maximum--Minimum Distance (MMD) descriptor was introduced for
human sleep staging by \citet{e18090272}; here we evaluate it for rodent
vigilance staging. It summarizes within-epoch amplitude excursion using a
sliding sub-window, and is intended to separate the high-amplitude
synchronized oscillations of SWS from the lower-amplitude, more variable
waveforms of REM and wake without reference to any particular frequency
band.

Let $W$ be the sub-window length in samples and $s$ the stride, giving
$K_{\mathrm{MMD}} = \lfloor (N - W)/s \rfloor + 1$ sub-windows
$\mathcal{W}_k = \{(k-1)s + 1, \dots, (k-1)s + W\}$. Within sub-window $k$
we locate the extremal samples and their positions,
%
\begin{equation}
\label{eq:mmdargs}
n_k^{+} = \operatorname*{arg\,max}_{n \in \mathcal{W}_k} x_n,
\qquad
n_k^{-} = \operatorname*{arg\,min}_{n \in \mathcal{W}_k} x_n,
\end{equation}
%
and form a per-sub-window distance $d_k$. The epoch-level feature is an
aggregation of these distances over the sub-windows,
%
\begin{equation}
\label{eq:mmd}
\mathrm{MMD}(x) \;=\; \Phi\!\left( d_1, \dots, d_{K_{\mathrm{MMD}}} \right).
\end{equation}
%
The manuscript describes $d_k$ as ``the statistical distance between
maximum and minimum signal values,'' which admits two instantiations that
give different numbers:
%
\begin{align}
d_k^{\mathrm{amp}} &= x_{n_k^{+}} - x_{n_k^{-}},
\label{eq:mmdamp}\\
d_k^{\mathrm{euc}} &= \sqrt{\,\alpha^{2}\big(n_k^{+} - n_k^{-}\big)^{2}
                     + \big(x_{n_k^{+}} - x_{n_k^{-}}\big)^{2}\,}.
\label{eq:mmdeuc}
\end{align}
%
Equation~(\ref{eq:mmdamp}) is the peak-to-peak amplitude range within the
sub-window. Equation~(\ref{eq:mmdeuc}) is the Euclidean distance between
the two extremal points in the time--amplitude plane, where $\alpha$ is the
scaling that maps a sample offset onto the time axis ($\alpha = 1$ for
sample index, $\alpha = 1/f_s$ for seconds). The second form is not
invariant to $\alpha$ or to the amplitude units, so both must be stated for
the feature to be reproducible.

% REQUIRES_EXPERIMENT: transcribe from the feature-extraction code, and
% then delete the sentence below and state the values directly:
%   (1) the sub-window length W and stride s (are the sub-windows
%       overlapping or contiguous?);
%   (2) which of Eq. (9) or Eq. (10) is computed, and if Eq. (10), the
%       value of alpha and the amplitude units;
%   (3) the aggregation Phi -- sum, mean, or max over sub-windows.
%       Results (line 329) calls MMD "a single scalar per epoch", so some
%       aggregation is applied, but the manuscript never says which.
% The source work of Aboalayon et al. is reported in the secondary
% literature as using a 1 s sliding sub-window with a summed Euclidean
% distance. That is the LITERATURE convention, NOT a measurement of this
% pipeline. Do not copy it into the paper as though it were this paper's
% setting -- check the code.
\noindent
We use $W = [\;]$ samples ($[\;]$\,s) with stride $s = [\;]$, the distance
of Equation~([\;]), and $\Phi = [\;]$.


\subsection{Cross-Frequency Coupling}
\label{subsec:cfc}

% =====================================================================
% !! READ THIS BEFORE INCLUDING THIS SUBSECTION !!
%
% This subsection is a SLOT, not a claim. As of the submitted version,
% the Methodology defines no coupling quantity, and the model's own input
% enumeration in Figures 3 and 4 (complexity, mobility, alpha/delta/
% gamma/beta/theta_power, activity, animal, mmd, entropy, eeg_data, plus
% raw sample indices) contains no coupling measure of any kind.
%
% Three outcomes, per section 6.1 of the revision plan:
%
%  (a) The code DOES compute a coupling quantity that never reached the
%      figures. Then fill the slots below, regenerate the figures from
%      the model that contains it, and the abstract becomes accurate.
%      This subsection ships as written. -- TRACK A.
%
%  (b) No coupling quantity is computed and the authors compute one and
%      retrain. This subsection ships once that is done. -- TRACK B.
%
%  (c) No coupling quantity is computed and none is added. Then DELETE
%      this entire subsection, delete the CFC row from
%      Table~\ref{tab:features}, delete the CFC line from
%      Algorithm~\ref{alg:pipeline}, and delete the CFC claims at lines
%      115 and 129 of the manuscript. The frozen abstract will then
%      overstate the method by one clause, which cannot be repaired by
%      writing -- but a body that silently repeats an unsupported claim
%      is strictly worse than a body that quietly does not make it.
%
% Under NO circumstances should a coupling formula be written here that
% the pipeline does not compute.
% =====================================================================
Cross-frequency coupling (CFC) quantifies the dependence between activity
in a low-frequency \emph{phase} band $B_{\varphi}$ and a high-frequency
\emph{amplitude} band $B_{a}$ \citep{canolty2010}. The epoch is band-pass
filtered into the two bands, and the analytic signal obtained by the
Hilbert transform $\mathcal{H}$ supplies an instantaneous phase and an
instantaneous amplitude envelope,
%
\begin{equation}
\label{eq:analytic}
\varphi(t) = \arg \mathcal{H}\big[x_{B_{\varphi}}\big](t),
\qquad
a(t) = \big| \mathcal{H}\big[x_{B_{a}}\big](t) \big|.
\end{equation}
%
A coupling statistic is then a scalar functional of the joint distribution
of $\varphi$ and $a$ over the epoch. The two standard choices are the
modulation index of \citet{tort2010}, which bins the phase into $M$ equal
bins, forms the mean envelope $\bar{a}_m$ in each bin, normalizes it to a
distribution $q_m = \bar{a}_m / \sum_{m'} \bar{a}_{m'}$, and measures its
divergence from uniformity,
%
\begin{equation}
\label{eq:mi}
\mathrm{MI} \;=\; \frac{\log M + \sum_{m=1}^{M} q_m \log q_m}{\log M},
\end{equation}
%
and the mean vector length,
%
\begin{equation}
\label{eq:mvl}
\mathrm{MVL} \;=\;
\Big| \tfrac{1}{N} \textstyle\sum_{t=1}^{N} a(t)\, e^{\mathrm{j}\varphi(t)} \Big|.
\end{equation}
%
% REQUIRES_EXPERIMENT: this is prerequisite P1 of the revision plan and
% must be answered from the feature-extraction code before this subsection
% is included at all. Required to fill the sentence below:
%   (1) which statistic is computed -- Eq. (12), Eq. (13), or another;
%   (2) the band pair(s) (B_phi, B_a) and how many pairs, since this
%       determines the per-epoch dimensionality of the CFC block;
%   (3) the band-pass filter design (type, order, zero-phase or not) and,
%       for Eq. (12), the phase-bin count M;
%   (4) whether any surrogate-based normalization was applied.
% If the answer to (1) is "none", see the guard comment above and delete
% this subsection rather than filling it.
\noindent
We compute [statistic] over the band pair(s) [$B_{\varphi}$, $B_{a}$],
yielding [$k$] coupling features per epoch.


% =====================================================================
% FEATURE TABLE
% This is the "one row per column of the design matrix" table of item
% A6.6. Dagger-marked cells are the ones that cannot be filled without
% reading the extraction code; the caption says so explicitly rather than
% carrying an invented number.
%
% If this becomes Table 1, renumber the existing performance table in the
% Results section and update its \ref.
% =====================================================================
\begin{table}[t]
\centering
\footnotesize
\setlength{\tabcolsep}{4pt}
\begin{tabular}{@{}lcp{3.0cm}@{}}
\toprule
\textbf{Descriptor} & \textbf{Dim.} & \textbf{Intended content} \\
\midrule
\multicolumn{3}{@{}l}{\emph{Spectral, from the Welch PSD, Eq.~(\ref{eq:welch})}}\\
Absolute band power $P_b$   & 5            & Oscillatory energy per canonical band \\
Relative band power $\tilde{P}_b$ & $\dagger$ & Spectral shape, invariant to overall gain \\
Peak frequency $f_b^{\star}$ & $\dagger$   & Dominant rhythm within a band \\
Spectral entropy $H$        & $\dagger$    & Broadband vs.\ rhythmic character \\
\addlinespace
\multicolumn{3}{@{}l}{\emph{Time-domain}}\\
Hjorth $A, M, C$            & 3            & Variance, normalized mean frequency, spectral bandwidth \\
$\mathrm{MMD}$              & 1            & Within-epoch amplitude excursion \\
\addlinespace
\multicolumn{3}{@{}l}{\emph{Coupling}}\\
CFC statistic               & $\dagger$    & Phase--amplitude dependence across bands \\
\addlinespace
\multicolumn{3}{@{}l}{\emph{Non-signal columns present in the fitted model}}\\
Animal index                & $\dagger$    & Categorical identifier of the recorded subject \\
Raw-sample columns          & $\dagger$    & Untransformed epoch samples \\
\midrule
\textbf{Total} $d$          & $\dagger$    & \\
\bottomrule
\end{tabular}
\caption{Components of the feature map $\phi$. Entries marked $\dagger$
are determined by implementation choices that the present description does
not fix and are stated in the released extraction code. The final two
blocks are not signal descriptors: they are additional columns of the
design matrix, disclosed here because they are visible in the attribution
analysis reported later.}
\label{tab:features}
\end{table}
% REQUIRES_EXPERIMENT: fill every dagger from the design matrix that
% produced the reported model, and make the total agree with the "5,009
% features" figure quoted for the network input. Note that the arithmetic
% does not currently close under any reading: 5,000 raw samples plus the
% 11 named columns visible in the SHAP figure is 5,011, not 5,009. If the
% classifier and the neural-network baseline did not receive the same
% columns, that must be stated here and the model comparison reframed
% accordingly.


% =====================================================================
% ALGORITHM BOX (item A6.7)
% Requires \usepackage{algorithm} and \usepackage{algorithmic}; see the
% notice at the top of this file.
% =====================================================================
\begin{algorithm}[t]
\caption{Vigilance-state classification pipeline}
\label{alg:pipeline}
\begin{algorithmic}[1]
\REQUIRE Single-channel EEG recording $r$ from animal $a$; sampling rate
         $f_s = 500$\,Hz; epoch length $T = 10$\,s
\ENSURE  Per-epoch posterior $\hat{p} \in \Delta^{2}$ and label
         $\hat{y} \in \mathcal{Y}$
\STATE $r \leftarrow \textsc{Preprocess}(r)$
       \COMMENT{filter cascade, notch, artifact rule}
\STATE Segment $r$ into non-overlapping epochs
       $x^{(1)}, \dots, x^{(E)} \in \mathbb{R}^{N}$, $N = T f_s$
\FOR{each epoch $x = x^{(e)}$}
  \STATE $\hat{S} \leftarrow \textsc{Welch}(x; w, L, D, N_{\mathrm{FFT}})$
         \COMMENT{Eq.~(\ref{eq:welch})}
  \FOR{each band $b \in \{\delta, \theta, \alpha, \beta, \gamma\}$}
    \STATE $P_b, \tilde{P}_b, f_b^{\star} \leftarrow$
           Eqs.~(\ref{eq:bandpower})--(\ref{eq:peakfreq})
  \ENDFOR
  \STATE $H \leftarrow$ Eq.~(\ref{eq:spectralentropy})
  \STATE $A, M, C \leftarrow$ Eqs.~(\ref{eq:activity})--(\ref{eq:complexity})
  \STATE $\mathrm{MMD} \leftarrow$ Eq.~(\ref{eq:mmd})
  \STATE $\mathrm{CFC} \leftarrow$ Eq.~(\ref{eq:mi}) or (\ref{eq:mvl})
         \COMMENT{omit if not computed}
  \STATE $\phi(x) \leftarrow
          \big[\{P_b, \tilde{P}_b, f_b^{\star}\}_b, H, A, M, C,
          \mathrm{MMD}, \mathrm{CFC}\big]$
\ENDFOR
\STATE Partition epochs into training and test sets
       \COMMENT{protocol stated in the experimental setup}
\STATE Fit gradient-boosted trees \citep{Chen_2016} on the training
       features and labels
\STATE $\hat{p} \leftarrow f(\phi(x))$;\;
       $\hat{y} \leftarrow \operatorname*{arg\,max}_{y \in \mathcal{Y}} \hat{p}_y$
\end{algorithmic}
\end{algorithm}
% REQUIRES_EXPERIMENT: line 1 (Preprocess) is written as a named step
% because the manuscript describes no preprocessing at all -- the words
% filter, notch, detrend, re-reference and artifact rejection appear
% nowhere, and 8 animals x 48 h x 360 epochs/h = 138,240 exactly, meaning
% no epoch was excluded. Either specify the cascade (type, cutoffs, order,
% zero-phase, mains notch frequency, reference scheme, rejection rule and
% the number of epochs excluded), or replace line 1 with an explicit
% statement that the signal was used unprocessed and add the corresponding
% limitation. Do not leave the step nameless.
% Line 15 (Partition) deliberately does not name a split: the manuscript
% currently describes four mutually inconsistent protocols. Name exactly
% one here once item A7 of the revision plan is done.
